\documentclass[11pt]{article}

\usepackage[margin=1in]{geometry}
\usepackage[T1]{fontenc}
\usepackage{amsmath}
\usepackage{booktabs}
\usepackage{hyperref}

\title{Research Report: Polynomial Ridge Regression on a Fixed Training and Validation Split}
\author{Agent Laboratory}
\date{}

\begin{document}

\maketitle

\section{Abstract}
Selecting polynomial complexity and regularization strength requires balancing predictive flexibility against estimation error while avoiding unsupported conclusions from limited validation evidence. We study this problem on a fixed public regression task using polynomial ridge models, with degree \(d\in\{1,\ldots,8\}\), regularization strength \(\alpha\in[0,100]\), and validation mean squared error (MSE) as the model-selection objective. The reported experiment evaluates a degree-one model with \(\alpha=0.1\), fitted on 80 training observations and assessed on 64 validation observations. The host execution reports training MSE \(0.00518281274299233\) and validation MSE \(0.007427150027746186\), with independent verification marked valid. The validation--training MSE difference is approximately \(0.00224434\); this descriptive gap alone does not establish overfitting. Our contribution is a preliminary empirical characterization of this configuration together with an explicit separation of measured performance, methodological motivation, and untested hypotheses. The supplied ridge-regression literature motivates examining coefficient shrinkage and held-out error, but does not establish an optimal degree or regularization strength for this task. Although the host records two actual CPU attempts, only one distinct configuration has been evaluated, so the available evidence supports neither comparative superiority nor independent replication. No confidence intervals, significance tests, or hidden-test measurements are available. We propose subsequent host-run comparisons within the permitted configuration space to assess whether additional polynomial terms improve validation performance, while accounting for the potential optimism introduced by repeated validation-based selection. The present result provides a verified reference measurement for that investigation, rather than a demonstrated model-selection optimum.

\section{Introduction}
Polynomial regression provides a direct setting for studying how model complexity and coefficient regularization affect prediction on observations excluded from fitting. Increasing polynomial degree expands the class of functions a model can represent, while ridge regularization penalizes coefficient magnitude. These choices can affect both approximation error and sensitivity to the available training sample. Greater flexibility may improve prediction when the response depends nonlinearly on the input, but it can also introduce unnecessary parameters when a simpler relationship is adequate. Consequently, training error alone is insufficient for choosing polynomial degree or regularization strength. This paper considers the fixed public regression task \texttt{dev-linear}, version \texttt{polynomial-ridge-1}, with seed \(11\). The task permits degrees \(d\in\{1,\ldots,8\}\) and regularization strengths \(\alpha\in[0,100]\), and defines validation mean squared error (MSE) as the selection objective. Our research question is whether increasing polynomial complexity or changing regularization improves validation performance relative to the evaluated degree-one configuration. The available experiment establishes a reference measurement for this question; it does not complete the comparison across the permitted model family.

For a scalar input \(x\), a polynomial predictor can be expressed as \(f_{d,\alpha}(x)=b+\sum_{j=1}^{d}\beta_j x^j\), where the coefficients are estimated from training observations under the host's ridge fitting procedure. The degree controls the available polynomial terms, and \(\alpha\) controls the strength of the quadratic coefficient penalty. Their effects depend on the observed inputs, feature covariance, and response variation, so an appropriate value cannot be inferred from the configuration parameters alone. The supplied synopsis of \emph{Ridge Regression: Structure, Cross-Validation, and Sketching} motivates examining regularization through held-out or cross-validation error and explains the dependence of ridge estimation on feature covariance and penalty strength (arXiv 1910.02373v2). That literature provides methodological context rather than task-specific evidence for an optimal configuration. For this task, a candidate's validation objective is
\[
\operatorname{MSE}_{\mathrm{val}}(d,\alpha)
=
\frac{1}{64}\sum_{i=1}^{64}
\left(y_i^{\mathrm{val}}-
\widehat f_{d,\alpha}(x_i^{\mathrm{val}})\right)^2,
\]
where \(\widehat f_{d,\alpha}\) is fitted using the training split. Comparing candidates requires actual host evaluations of this quantity. Repeatedly selecting candidates against the same validation observations also creates a distinction between the lowest observed validation error and performance on a separate, unused sample.

The reported experiment uses \(80\) training observations and \(64\) validation observations with the configuration \(d=1\) and \(\alpha=0.1\). The supplied host record identifies the run as \texttt{execution-0002}, reports actual CPU execution, and marks independent verification as valid. Its training MSE is \(0.00518281274299233\), and its validation MSE is \(0.007427150027746186\). The validation--training difference is approximately \(0.00224434\). These measurements characterize the fitted configuration on the supplied split. Although validation error exceeds training error, this difference alone does not establish overfitting: the two quantities are computed on different observations, and the training observations were used to estimate the model. The host reports two actual CPU attempts and one distinct configuration, but those counts do not establish independent replication or a comparison between alternative models. No confidence interval, standard error, or significance test accompanies the aggregate metrics. Hidden test observations remain withheld, and no test performance is reported.

The paper makes three limited contributions. First, it documents the evaluated configuration and its verified training and validation errors on a fixed public split. Second, it connects polynomial complexity and coefficient shrinkage to the task's validation objective while distinguishing literature-based motivation from local experimental findings. Third, it specifies a prospective comparison procedure in which the host evaluates additional permitted configurations, retains their measured results, and selects the lowest observed validation MSE among the evaluated candidates. A hypothesis motivating those comparisons is that a degree-one model provides sufficient predictive capacity and that additional polynomial terms may primarily capture sample variation. This hypothesis remains untested. Comparisons across degrees and regularization strengths would be needed to assess it, and any resulting selection would remain conditional on the candidates examined and the fixed validation sample. The present study therefore supplies a preliminary empirical reference and a clearly bounded model-selection question. It does not establish an optimal degree, an optimal regularization strength, an improvement over a measured baseline, or generalization to the withheld test split.

\section{Background}
We consider supervised regression with scalar inputs \(x\in\mathbb{R}\) and responses \(y\in\mathbb{R}\). Let \(\mathcal{D}_{\mathrm{tr}}=\{(x_i^{\mathrm{tr}},y_i^{\mathrm{tr}})\}_{i=1}^{80}\) denote the supplied training observations and \(\mathcal{D}_{\mathrm{val}}=\{(x_i^{\mathrm{val}},y_i^{\mathrm{val}})\}_{i=1}^{64}\) denote the fixed validation observations. These splits belong to task \texttt{dev-linear}, version \texttt{polynomial-ridge-1}, with seed \(11\). Polynomial regression represents an input through the feature vector \(\phi_d(x)=(x,x^2,\ldots,x^d)^\top\), yielding the predictor
\[
f_{d}(x;b,\boldsymbol{\beta})
=
b+\phi_d(x)^\top\boldsymbol{\beta},
\qquad
\boldsymbol{\beta}\in\mathbb{R}^{d}.
\]
The permitted degree satisfies \(d\in\{1,\ldots,8\}\). A degree-one predictor is affine in the input, whereas higher degrees introduce curvature through additional powers of \(x\). Although these predictors can be nonlinear functions of the input, their dependence on the fitted coefficients remains linear. The corresponding unregularized function classes are nested: every degree-\(d\) polynomial is also representable at degree \(d+1\) by setting the additional coefficient to zero. Consequently, increasing degree cannot increase the minimum attainable unpenalized training squared error when the same intercept convention is used and optimization is exact. This representational property does not imply decreasing validation error. Additional coefficients must be estimated from a finite training sample, and their estimation can introduce variation without providing useful predictive structure. The supplied observations motivate examining a simple affine model, but neither the task identifier nor qualitative inspection establishes that the underlying response function is linear. No particular noise distribution or known data-generating regression function is assumed in the empirical interpretation.

Ridge regression supplements squared-error fitting with a quadratic penalty on coefficients. Under a conventional formulation with an unpenalized intercept and a summed training loss, the estimator is
\[
(\widehat b,\widehat{\boldsymbol{\beta}})
\in
\operatorname*{arg\,min}_{b,\boldsymbol{\beta}}
\left\{
\sum_{i=1}^{80}
\left(y_i^{\mathrm{tr}}-b-\phi_d(x_i^{\mathrm{tr}})^\top
\boldsymbol{\beta}\right)^2
+
\lambda\|\boldsymbol{\beta}\|_2^2
\right\}.
\]
Here \(\lambda\geq0\) denotes the penalty in this mathematical convention. The experimental interface instead specifies the host parameter \(\alpha\in[0,100]\); its exact correspondence to \(\lambda\) depends on the host's loss normalization and intercept treatment. For example, replacing the summed loss with its average changes the relative penalty strength unless the penalty parameter is rescaled accordingly. These conventions are therefore relevant to reproducing a numerical setting such as \(\alpha=0.1\). The displayed objective introduces the standard ridge principle and does not independently establish the host's implementation details. For centered training features and responses under this conventional objective, let \(Z_d\in\mathbb{R}^{80\times d}\) be the centered feature matrix and \(\boldsymbol{y}_c\in\mathbb{R}^{80}\) the centered response vector. When \(\lambda>0\), the coefficient solution is
\[
\widehat{\boldsymbol{\beta}}
=
\left(Z_d^\top Z_d+\lambda I_d\right)^{-1}
Z_d^\top\boldsymbol{y}_c.
\]
The penalty makes this coefficient system invertible even when feature columns are linearly dependent. At \(\lambda=0\), the objective reduces to ordinary least squares, for which coefficient uniqueness requires appropriate rank conditions. These statements describe estimator properties rather than measurements from the reported execution.

The effect of regularization depends on the geometry of the feature representation. Writing a singular value decomposition as \(Z_d=U S V^\top\), ridge estimation weights the response along each represented feature direction by \(s_k/(s_k^2+\lambda)\), where \(s_k\) is the corresponding singular value. Equivalently, the fitted centered responses have the form
\[
\widehat{\boldsymbol{y}}_c
=
H_\lambda\boldsymbol{y}_c,
\qquad
H_\lambda
=
Z_d\left(Z_d^\top Z_d+\lambda I_d\right)^{-1}Z_d^\top,
\qquad
\operatorname{tr}(H_\lambda)
=
\sum_k\frac{s_k^2}{s_k^2+\lambda}.
\]
The trace gives a conventional measure of the effective degrees of freedom of the penalized feature fit, excluding the unpenalized intercept. Increasing the penalty reduces this quantity for a fixed feature matrix. Increasing polynomial degree changes the matrix itself, so complexity cannot be characterized by the penalty parameter alone. Polynomial columns can also be strongly correlated, particularly when successive powers have similar behavior over the observed input range. Ridge penalization can reduce sensitivity along poorly determined directions, while simultaneously introducing shrinkage bias. Moreover, the Euclidean coefficient penalty is sensitive to feature scaling: rescaling a polynomial column changes the penalty associated with an equivalent prediction function. These considerations explain why degree and regularization should be studied jointly under a fixed fitting procedure. The supplied synopsis of \emph{Ridge Regression: Structure, Cross-Validation, and Sketching} provides relevant context for covariance-dependent ridge estimation and validation-based regularization selection (arXiv 1910.02373v2). Its discussion of high-dimensional validation bias does not supply a quantitative prediction for this small polynomial task, and no singular values, effective degrees of freedom, or conditioning measurements are available from the supplied aggregate execution metrics.

Predictive assessment uses mean squared error on observations excluded from coefficient fitting. For a candidate fitted on \(\mathcal{D}_{\mathrm{tr}}\), define
\[
\widehat R_{\mathrm{val}}(d,\alpha)
=
\frac{1}{64}
\sum_{i=1}^{64}
\left(
y_i^{\mathrm{val}}
-
\widehat f_{d,\alpha}(x_i^{\mathrm{val}})
\right)^2.
\]
This empirical quantity is distinct from the population risk \(R(\widehat f)=\mathbb{E}_{(X,Y)\sim P}[(Y-\widehat f(X))^2]\), where \(P\) denotes an underlying distribution if such a sampling model is adopted. Relating the two requires assumptions about sampling and distributional agreement that are not established by the aggregate metrics. Training MSE is additionally affected by using the same observations to estimate coefficients; a positive validation--training gap therefore does not, by itself, diagnose excessive model complexity. If a finite candidate set \(\mathcal C\) is evaluated, validation-based selection chooses a member minimizing \(\widehat R_{\mathrm{val}}\) within that set. Such a choice need not minimize population risk or attain the best validation error over every permitted degree and real-valued regularization strength. Repeated adaptive use of the same validation split can also favor candidates whose errors happen to be low on those particular observations. The present record contains only one distinct evaluated configuration, \(d=1\) and \(\alpha=0.1\), so it supplies no empirical ordering of alternative candidates. The withheld test split contains \(96\) observations according to the manifest, but its predictive error is unmeasured. These distinctions define the scope of the study: the verified execution characterizes a fitted candidate on one fixed training and validation split, while comparisons across configurations and conclusions about independent generalization require additional evidence.

\section{Related Work}
\emph{Ridge Regression: Structure, Cross-Validation, and Sketching} provides the closest methodological context for the present study (arXiv 1910.02373v2). The supplied synopsis describes squared-error regression with a quadratic coefficient penalty, the dependence of ridge estimation on feature covariance and regularization strength, and the selection of regularization through held-out or cross-validation error. Our study uses the same general fitting principle but addresses a restricted empirical question: how polynomial degree \(d\in\{1,\ldots,8\}\) and ridge strength \(\alpha\in[0,100]\) affect validation MSE on a fixed public split. Increasing degree changes the feature representation, whereas changing \(\alpha\) changes coefficient shrinkage within that representation. Consequently, the two parameters should be considered jointly rather than treating regularization as independent of model complexity. The literature motivates this comparison but supplies no optimal configuration for the present observations. Our available measurement concerns only \(d=1\) and \(\alpha=0.1\); it therefore documents one member of the permitted family rather than empirically validating a general model-selection rule.

The distinction between validation protocols also limits the comparison with that work. The supplied synopsis discusses validation bias in a high-dimensional asymptotic regime, while our experiment uses \(80\) training observations and \(64\) validation observations with a single scalar input expanded into polynomial features. We do not establish that the assumptions of the cited asymptotic analysis hold in this setting, so its quantitative implications cannot be transferred directly. Moreover, our reported objective comes from one fixed held-out split, not an executed cross-validation procedure. The applicable methodological lesson is to evaluate prediction on observations excluded from coefficient fitting and to distinguish selection error from independent generalization evidence. A prospective comparison would retain the host's fitting procedure and split while varying degree and regularization, selecting
\[
(\widehat d,\widehat\alpha)
\in
\operatorname*{arg\,min}_{(d,\alpha)\in\mathcal C}
\operatorname{MSE}_{\mathrm{val}}(d,\alpha),
\]
where \(\mathcal C\) contains only configurations actually evaluated by the host. This expression defines a proposed selection procedure, not a completed search or an optimum over the entire permitted space. Sketching is another topic covered by the cited study, but the supplied local experiment contains no sketching implementation or approximation comparison. Accordingly, we make no claim about sketching accuracy or computational benefit.

\emph{Agent Laboratory: Using LLM Agents as Research Assistants} addresses the organization of research rather than the statistical choice of a regression model (arXiv 2501.04227v2). According to the supplied summary, specialized agents perform literature review, planning, code experimentation, interpretation, and reporting. The present task shares that separation of activities, but its execution interface is substantially narrower: candidate experiments are specified by literal degree and regularization configurations, and the trusted host performs fitting and validation. This division assigns different evidential roles to a proposal, a literature statement, and an execution receipt. For example, the reported validation MSE \(0.007427150027746186\) is a host measurement for the evaluated configuration; the suggestion that additional polynomial terms may fit sample variation is an untested hypothesis. The cited workflow is relevant to maintaining this distinction, but it is not a predictive baseline for polynomial ridge regression. No controlled comparison of agent workflows, reporting accuracy, or research productivity was performed here.

\emph{AgentRxiv: Towards Collaborative Autonomous Research} concerns sharing and retrieving prior research reports through a local preprint server and similarity-based retrieval, according to the supplied summary (arXiv 2503.18102v1). Its relationship to our study lies in access to prior research evidence rather than an alternative estimator. Retrieving a related report can inform candidate selection or interpretation, but the applicability of its findings still depends on the data, feature representation, fitting convention, and evaluation protocol. Our review uses the supplied frozen literature descriptions and does not evaluate a report-sharing system or retrieval method. Thus, neither AgentRxiv nor Agent Laboratory supports a claim that the measured ridge configuration outperforms another model. Together, these workflow references motivate explicit provenance and separation of research stages, while the ridge reference motivates the statistical comparison itself. The remaining empirical question is whether alternative permitted configurations improve validation MSE under the same host protocol. Answering it requires additional measured candidates; the current literature and execution record do not establish comparative superiority, statistical significance, or withheld-test performance.

\section{Methods}
We study polynomial ridge regression using the fixed training and validation splits supplied for task \texttt{dev-linear}, version \texttt{polynomial-ridge-1}, with seed \(11\). The training split contains \(80\) observations and the validation split contains \(64\) observations. For a candidate configuration \(c=(d,\alpha)\), the admissible parameter space is
\[
\mathcal{H}
=
\{1,\ldots,8\}\times[0,100],
\qquad
\phi_d(x)=(x,x^2,\ldots,x^d)^\top.
\]
We write the fitted predictor as
\[
\widehat f_c(x)
=
\widehat b_c+\phi_d(x)^\top\widehat{\boldsymbol{\beta}}_c,
\qquad
(\widehat b_c,\widehat{\boldsymbol{\beta}}_c)
=
\operatorname{Fit}_{\mathrm{host}}(\mathcal{D}_{\mathrm{tr}};d,\alpha).
\]
Here \(\operatorname{Fit}_{\mathrm{host}}\) denotes the trusted host's polynomial ridge fitting procedure. This notation separates the experimental interface from the conventional ridge objective introduced in Background. In particular, the supplied execution summary does not establish the host's loss normalization, intercept penalization, feature scaling, or numerical solver. We therefore retain \(\alpha\) as the host parameter rather than asserting an unverified correspondence with the theoretical penalty \(\lambda\). Coefficients are fitted using the training observations, and validation observations provide the prescribed prediction assessment. The evaluated configuration is \(c_0=(1,0.1)\), specified by the literal assignment \texttt{CONFIG = \{'degree': 1, 'alpha': 0.1\}}. Its predictor has the affine form \(\widehat f_{c_0}(x)=\widehat b_{c_0}+\widehat\beta_{c_0}x\). This model provides a simple reference for studying whether additional polynomial terms offer useful predictive capacity; its use does not assume that the unknown response function is linear.

For each host-evaluated candidate, assessment uses the unpenalized mean squared prediction error on each supplied split. The coefficient penalty belongs to the fitting procedure and is not added to either reported MSE. With \(n_{\mathrm{tr}}=80\) and \(n_{\mathrm{val}}=64\), define
\[
e_{i,s}(c)
=
y_i^s-\widehat f_c(x_i^s),
\qquad
\widehat R_s(c)
=
\frac{1}{n_s}\sum_{i=1}^{n_s}e_{i,s}(c)^2,
\qquad
s\in\{\mathrm{tr},\mathrm{val}\}.
\]
The selection objective is \(\widehat R_{\mathrm{val}}\); training MSE supplies a descriptive assessment of fit on observations used to estimate coefficients. We also define the descriptive difference
\[
\Delta(c)=\widehat R_{\mathrm{val}}(c)-\widehat R_{\mathrm{tr}}(c).
\]
This difference compares errors on two separate samples and does not isolate an effect of excessive complexity. Neither its sign nor its magnitude supplies a significance test. The reported evidence consists of aggregate errors, so the analysis does not estimate residual distributions, observation-level uncertainty, confidence intervals, or statistical significance. No cross-validation, resampling, or alternative splitting procedure is reported as executed. The supplied ridge-regression synopsis motivates evaluating regularization through held-out prediction error (arXiv 1910.02373v2), but the local protocol remains a single fixed validation split.

Comparative model selection is prospective because the host reports only one distinct evaluated configuration. To make a subsequent comparison concrete, we propose the finite candidate set
\[
\mathcal{C}_{\mathrm{proposed}}
=
\{1,\ldots,8\}\times
\{0,0.01,0.1,1,10,100\}.
\]
This set contains \(48\) configurations, including \(c_0\), and would require evaluating \(47\) additional distinct configurations if the existing verified measurement is retained. That proposal is numerically compatible with the reported \(49\) remaining CPU executions, although actual scheduling remains subject to the host's execution and model-request limits. No additional candidate is evaluated by this report. The proposed design varies degree and regularization jointly while retaining the same data and host fitting procedure. Including \(\alpha=0\) provides the unregularized endpoint, while positive values sample several penalty scales. The grid is a bounded comparison design rather than an exhaustive search over the continuous interval. For the subset \(\mathcal{C}_{\mathrm{measured}}\) actually evaluated successfully, selection would use
\[
\widehat c
\in
\operatorname*{arg\,min}_{c\in\mathcal{C}_{\mathrm{measured}}}
\widehat R_{\mathrm{val}}(c).
\]
As a prospective deterministic tie rule, exactly equal host-reported validation errors would favor the smaller degree and then the smaller \(\alpha\). This rule would resolve reporting ambiguity without establishing a statistically meaningful difference between tied candidates. The current singleton set \(\mathcal{C}_{\mathrm{measured}}=\{c_0\}\) supports retaining the evaluated configuration but provides no empirical comparison with alternatives.

Experimental provenance is maintained through the supplied host receipt rather than inferred from the proposed configuration or literature. The selected record identifies \texttt{execution-0002} as an actual CPU execution and marks independent verification as \texttt{valid}. Its reported training and validation MSE values are \(0.00518281274299233\) and \(0.007427150027746186\), respectively. These numbers are attributed to that named invocation. The separate host observations of two actual CPU attempts and one distinct configuration are retained as resource accounting; they are not treated as two independent statistical replications. For future comparisons, each candidate should remain associated with its configuration, execution identifier, verification status, and reported errors, including candidates that do not improve the selection objective. Any selected configuration would be described as the lowest observed validation-MSE candidate among those evaluated. Repeated selection against the same validation sample can introduce selection optimism, and the present protocol supplies no independent estimate of its magnitude. The manifest identifies \(96\) withheld test observations, but test access remains withheld and no test error enters fitting, selection, or reporting. Thus, the methodology distinguishes the completed host measurement from the proposed comparison and preserves the limits of conclusions supported by the available evidence.

\section{Experimental Setup}
The experiment uses the fixed public regression task \texttt{dev-linear}, task version \texttt{polynomial-ridge-1}, with seed \(11\). Each supplied observation consists of a scalar input \(x\), a scalar response \(y\), and an identifier recording its task, seed, split, and observation index. The host supplies \(80\) training observations and \(64\) validation observations; the split manifest additionally records \(96\) test observations whose values and predictive errors remain withheld. The manifest describes generation using the standard-library \texttt{random.Random} generator, with split order training, validation, and test. This information specifies the recorded generation protocol but does not establish a response-generating equation, a noise distribution, or independence assumptions sufficient for statistical inference. We use the supplied split assignments without proposing changes to their membership. Training observations provide the data for coefficient estimation, and validation observations provide the prescribed model-selection objective. Observation identifiers distinguish the records associated with each split, while the manifest supplies SHA-256 digests for the split contents and their identifiers. These digests provide recorded data provenance; their presence is distinct from independently recomputing them, which was not performed for this report. No external dataset, additional response labels, or hidden-test information is introduced. The resulting empirical setting is a single fixed training and validation partition, rather than a collection of independently generated datasets or repeated random partitions. Consequently, the experimental unit for the reported predictive assessment is the named host execution on these supplied observations. The setup does not treat the task name as evidence that an affine predictor is correctly specified.

The evaluated model is polynomial ridge regression with degree \(d=1\) and host regularization parameter \(\alpha=0.1\). These settings reproduce the original research plan and the configuration recorded for \texttt{execution-0002}. The admissible interface permits an integer degree from \(1\) through \(8\) and a finite regularization parameter from \(0\) through \(100\). The completed experiment therefore evaluates one permitted member of the model family, with predictor representation
\[
\widehat f_{1,0.1}(x)=\widehat b+\widehat\beta x.
\]
The configuration is submitted as the literal assignment \texttt{CONFIG = \{'degree': 1, 'alpha': 0.1\}}, and the trusted host performs polynomial feature construction, ridge fitting, and prediction assessment using the preloaded data. The proposal interface does not supply custom preprocessing code, a replacement solver, or an alternative training objective. Although the model can be represented by an intercept and one slope coefficient, the execution summary does not report their fitted values. It also does not establish feature scaling, intercept penalization, loss normalization, floating-point precision, software dependency versions, or the particular numerical linear algebra routine. We therefore describe these as unspecified implementation details rather than assigning conventional defaults to the host. In particular, the numerical parameter \(\alpha=0.1\) is interpreted within the host implementation; it is not equated without evidence to a penalty coefficient under another loss-normalization convention. The receipt identifies actual CPU execution, but processor model, memory allocation, operating-system configuration, and runtime are not supplied as experimental measurements. Reproduction of the reported result consequently requires the registered host task and data, in addition to the two configuration fields. The supplied artifact listing identifies source files and a registered specification that could support a later implementation audit, but this report does not claim to have inspected those artifacts.

Evaluation uses mean squared prediction error, calculated separately over the training and validation observations. For split \(s\), let \(n_s\) denote its recorded observation count and let \(\widehat f\) denote the predictor fitted on the training split. The assessment quantity is
\[
\operatorname{MSE}_{s}
=
\frac{1}{n_s}\sum_{i=1}^{n_s}
\bigl(y_i^{s}-\widehat f(x_i^{s})\bigr)^2,
\qquad
n_{\mathrm{tr}}=80,\quad n_{\mathrm{val}}=64.
\]
Each observation receives equal weight in this definition, and the reported metric is a mean rather than a total squared error. The regularization penalty used during fitting is not included in the predictive MSE. Training MSE describes performance on observations used to estimate coefficients, whereas validation MSE is the designated objective for comparing permitted configurations. Both quantities have units of squared response values; their magnitude alone does not establish practical adequacy because no application-specific error tolerance is supplied. For traceability, the host reports training MSE \(0.00518281274299233\) and validation MSE \(0.007427150027746186\) for \texttt{execution-0002}, with independent verification marked \texttt{valid}. These values are execution evidence, rather than predictions derived from the literature or the proposed model family. Their descriptive difference is approximately \(0.00224434\). The setup provides no executed confidence-interval procedure, significance test, cross-validation run, or residual-based uncertainty analysis. Aggregate MSE values alone do not identify whether errors are concentrated in a few observations, vary across the input domain, or arise from systematic prediction bias. The artifact listing includes model and prediction files, but no additional diagnostics are inferred from their names. The verification status is reported as supplied and does not itself quantify uncertainty in population prediction error. A positive validation--training difference is therefore recorded descriptively without using it as a stand-alone diagnostic of overfitting.

The experimental accounting separates completed execution from the prospective comparison design described in Methods. The host records two actual CPU attempts, one distinct configuration, a limit of \(51\) actual CPU executions, and \(49\) remaining executions. The supplied ledger identifies receipts for \texttt{execution-0000} and \texttt{execution-0002}; the aggregate metrics reported here belong specifically to the latter invocation. These accounting observations do not establish independent statistical replication, different training samples, or measurements for competing hyperparameters. Cached replay is not a new CPU attempt under the host contract. A prospective comparison would use the previously specified grid \(\{1,\ldots,8\}\times\{0,0.01,0.1,1,10,100\}\), which contains \(48\) configurations. Retaining the existing measurement would leave \(47\) additional distinct configurations to evaluate, subject to the host's separate execution and model-request limits. This grid is an experimental proposal and has not been executed in the evidence supplied for this section. For comparable future assessments, the host task, split membership, fitting implementation, and metric definition would remain fixed while degree and regularization vary. Each successful candidate would be associated with its execution identifier, verification status, and both reported errors; unsuccessful evaluations would be retained as execution outcomes rather than assigned invented error values. Selection would minimize validation MSE among successfully measured candidates, with the prospective tie rule from Methods applied only to exactly equal reported values. Such selection would describe the best observed candidate within the evaluated set, without establishing a continuous-parameter optimum. Reusing the same validation sample would also limit its role as independent evidence after selection. The withheld test split supplies no measurement in the present setup, and neither the resource budget nor the proposed grid supports claims about computational speed, monetary cost, baseline improvement, or hidden-test accuracy.

\section{Results}
The host evaluated polynomial ridge regression with degree \(d=1\) and regularization parameter \(\alpha=0.1\) on the fixed \texttt{dev-linear} split with seed \(11\). The supplied record for \texttt{execution-0002} identifies an actual CPU execution and marks independent verification as \texttt{valid}. Table~\ref{tab:observed-errors} reports the supplied predictive errors. Coefficient fitting used \(80\) training observations, and validation assessment used \(64\) observations excluded from fitting. The reported values describe this configuration under the registered host procedure; they are not averages across datasets or independently replicated experiments.

\begin{table}[ht]
\centering
\caption{Host-reported mean squared errors for \texttt{execution-0002}, with \(d=1\) and \(\alpha=0.1\).}
\label{tab:observed-errors}
\begin{tabular}{lrr}
\toprule
Split & Observations & Mean squared error \\
\midrule
Training & 80 & \(0.00518281274299233\) \\
Validation & 64 & \(0.007427150027746186\) \\
\bottomrule
\end{tabular}
\end{table}

Validation MSE exceeds training MSE. Their descriptive difference, calculated from the supplied aggregate measurements, is
\[
\Delta
=
0.007427150027746186
-
0.00518281274299233
\approx 0.00224434.
\]
This difference compares prediction errors on separate observation sets, one of which was used to estimate the coefficients. It does not isolate the contribution of model complexity or establish overfitting. The supplied evidence contains no uncertainty estimates, confidence intervals, standard errors, or significance tests. Aggregate MSE also does not reveal whether individual errors are concentrated near particular input values or dominated by a small number of observations. Although prediction artifacts are listed in the receipt, their contents were not inspected for this report.

The host records two actual CPU attempts and one distinct configuration. The ledger identifies receipts for \texttt{execution-0000} and \texttt{execution-0002}, but the supplied numerical results reported here belong to the latter invocation. These accounting records do not establish independent statistical replication or comparative performance. No measured alternative degree, unregularized baseline, or alternative positive regularization strength is available. Consequently, the experiment does not provide an ablation of coefficient shrinkage, a comparison of polynomial capacities, or evidence that the evaluated hyperparameters minimize validation MSE across the permitted space. The proposed comparison grid remains unexecuted; its entries cannot be included as experimental results.

The observed validation error supplies a reference measurement for subsequent comparisons under the same fitting procedure, training observations, validation observations, and metric definition. Such controls would make candidate differences interpretable within this fixed task, although repeated validation selection could introduce optimism. The current evidence does not test the hypothesis that additional polynomial terms primarily fit sample variation. It also supplies no application-specific error threshold against which predictive adequacy could be judged. The \(96\) test observations recorded in the manifest remain withheld, and test performance is unmeasured. No runtime comparison, monetary cost, or baseline improvement is supported by the supplied execution evidence. The empirical conclusion is therefore limited to the verified training and validation errors of the degree-one, \(\alpha=0.1\) configuration.

\section{Discussion}
This study documents a preliminary evaluation of polynomial ridge regression on the fixed public task \texttt{dev-linear}, version \texttt{polynomial-ridge-1}, with seed \(11\). The evaluated configuration, \(d=1\) and \(\alpha=0.1\), produced training MSE \(0.00518281274299233\) on \(80\) observations and validation MSE \(0.007427150027746186\) on \(64\) observations. These measurements belong to the actual CPU invocation \texttt{execution-0002}, whose independent verification status is reported as \texttt{valid}. The validation--training difference is approximately \(0.00224434\). This difference describes performance on two observation sets with different roles in estimation; it does not establish excessive model complexity or identify its cause. The result provides a measured reference for subsequent comparisons, but the available evidence does not establish predictive adequacy relative to an application-specific tolerance. Retaining the evaluated configuration is justified as a record of completed evaluation, while choosing it as an optimal model would require comparative evidence.

The principal limitation is that only one distinct configuration has been evaluated. The host records two actual CPU attempts, but these counts do not demonstrate independent statistical replication or provide measurements for alternative hyperparameters. In particular, the experiment cannot separate the effect of using a degree-one representation from the effect of applying regularization at \(\alpha=0.1\). The hypothesis that an affine predictor supplies sufficient capacity, with additional polynomial terms primarily capturing sample variation, remains untested. The supplied ridge-regression literature motivates examining feature covariance, coefficient shrinkage, and held-out prediction error (arXiv 1910.02373v2), but it supplies no task-specific optimum. Its discussion of validation bias in a high-dimensional regime also does not quantify uncertainty for this experiment. No confidence interval, standard error, or significance test accompanies the reported aggregate errors, and independent verification of an execution is distinct from uncertainty assessment. The \(96\) test observations remain withheld, so population generalization and hidden-test performance cannot be inferred from the reported validation MSE.

A concrete next step is the prospective comparison over \(\{1,\ldots,8\}\times\{0,0.01,0.1,1,10,100\}\) described in Methods. This grid contains \(48\) configurations, including the evaluated configuration, and would require \(47\) additional distinct evaluations if its existing measurement is retained. That count fits within the reported \(49\) remaining CPU executions, subject to the host's separate request limits and successful execution of each candidate. These evaluations are proposed rather than completed. Holding the splits and fitting procedure fixed would allow comparisons across degrees at each penalty value and across penalties within each degree. Including \(\alpha=0\) would provide an unregularized comparison under the same host conventions. All outcomes should be retained, including candidates with higher error and unsuccessful evaluations. Selection should minimize validation MSE among successfully evaluated candidates, with the specified deterministic rule resolving exact ties. A resulting choice would establish the lowest observed validation error within this finite grid; it would not establish an optimum over the continuous regularization interval or a statistically significant advantage.

Further interpretation should account for repeated use of the validation split. Searching more candidates can favor configurations that happen to predict these particular validation observations well, even when their advantage would not persist on unused data. The present evidence does not measure the magnitude of this selection optimism. Under the current protocol, this limitation should remain explicit rather than being addressed through unauthorized resampling or access to withheld observations. A future study with separately authorized independent assessment could investigate whether any observed ordering persists beyond this split. Reproducibility would also benefit from documenting the host's feature scaling, intercept treatment, loss normalization, and solver, because these conventions determine the meaning of a numerical regularization parameter. Within the existing evidence, the contribution is a traceable measurement and a bounded comparison proposal. Establishing the value of additional polynomial capacity or coefficient shrinkage requires the proposed host-run evaluations; establishing independent generalization requires evidence beyond repeated selection on the same validation sample.

\end{document}